% Pursuing Stacks (cotes 134-2 to 134-8): archive cote 134-7, site page 82 (PDF page 83), typescript page 524.
% typescript. Draft compared: ps6.tex:4083-4170.
% AI-transcribed by Claude Opus 5.5 (claude-opus-5-5, adaptive thinking, effort high), Message Batches API, batch ps4 (msgbatch_012Aoc1NuskATPsyBeKSYK7J), 2026-11-12; not yet checked by eye.
\scanpage{134-7/82}
\mspage{524}
\[
(6\ b)\qquad \left\{\begin{array}{l}
\overset{\hand{\sim}}{S}(L\hand{\pi}.(X)) \hand{\longrightarrow} \ X \\
L. \hand{\to} L\hand{\pi}.(\overset{\hand{\sim}}{S}(L.)) \qquad .
\end{array}\right.
\]

The adjunction formulas for the two pairs in (5) are
\[
(7)\qquad \left\{\begin{array}{ll}
a) & Hom(L\overset{\hand{\sim}}{H}.(X),L.) \hand{\simeq} Hom(X, \overset{\hand{\sim}}{K}(L.)) \\
b) & Hom(L., L\hand{\pi}.(X)) \hand{\simeq} Hom(\overset{\hand{\sim}}{S}(L),X) \qquad .
\end{array}\right.
\]

The first hand side of a) should be viewed as \hand{\uline{cohomology}} of X with
coefficients\xout{x} in L. (more specifically, \inserted{as} $\overset{\hand{\sim}}{H}{}^{0}(X, L.)$, and in case $L. =\hand{\pi}\hand{[}n\hand{]}$\hand{,}
\xout{\illegible{1}} posing
\[
(8\ a)\qquad K(\hand{\pi}\hand{[}n\hand{]}) \overset{\text{def}}{=} K(\hand{\pi},n) \qquad ,
\]
we get that the Eilenberg-Maclane object $K(\hand{\pi},n)$ represents the cohomology
functor
\[
X \hand{\longmapsto} \overset{\hand{\sim}}{H}{}^{0}(X,\hand{\pi}\hand{[}n\hand{]}) = H^{n}(X,\hand{\pi}) \quad .
\]

Symetrically, the first hand side of (7 b) should be viewed as
a ``mixed homotopy \deleted{group}\inserted{module}'' of X (relative to the ``co-coefficient'' L.),
I am tempted to denote it \xout{\illegible{1}} as
\margin{\textcircled{H}}
\[
\hand{\pi}_{0}(X,L.)\ ,\quad resp.\ \ \hand{\pi}_{n}(X,\hand{H})\quad if\ \ L. =\hand{H}\hand{[}n\hand{]}\ ;
\]
thus, posing
\[
(8\ b)\qquad \overset{\hand{\sim}}{S}(\hand{H}\hand{[}n\hand{]}) \overset{\text{def}}{=} S(\hand{H},n) \qquad \text{(sph\inserted{e}re-like object for $\hand{H}$,n) ,}
\]
we get that the $\hand{H}$-sphere $S(\hand{H},n)$ \inserted{over k} represents the (covariant) functor \inserted{on $Hot_{0}(k)$}
\[
X \hand{\longmapsto} \hand{\pi}_{n}(X,\hand{H}) \quad .
\]
In case $\hand{H}=$\deleted{\illegible{1}}\hand{k}\hand{,} we get
\[
\hand{\pi}_{n}(X,k) \ = \ Hom(k\hand{[}n\hand{]},L\hand{\pi}.(X)) = H_{n}(L\hand{\pi}.(X)) \overset{\text{def}}{=} \hand{\pi}_{n}(X) \quad ,
\]
and we get that the \xout{\unsure{functor}} ``usual'' homotopy \deleted{group}\inserted{module} functor
\[
X \hand{\longrightarrow} \hand{\pi}_{n}(X)
\]
is represented by the ``usual'' n-sphere (over k however !) S(n,k), as
it should. ( This is of course the main justification why we expect
$(\overset{\hand{\sim}}{S},L\hand{\pi}.)$ to be a pair of adjoint functors, whereas adjunction for the
pair $(L\overset{\hand{\sim}}{H}.,\overset{\hand{\sim}}{K})$ is already fairly familiar from the set-up of ordinary
homotopy types.)

The adjunction formulae (7) show that the objects $\overset{\hand{\sim}}{K}(L.)$
\margin{\textcircled{$\widetilde{W}(X)$}}
in $H\unsure{o}t_{0}(k)$ a\unsure{r}e group objects, whereas the objects $\overset{\hand{\sim}}{S}(L.)$ are co-group
objects. This gives some \inserted{inner} justification for calling \inserted{$(\widetilde{W}(X)=$} $\overset{\hand{\sim}}{K}(L\overset{\hand{\sim}}{H}.(X))$ the
\margin{82}
